\exercisesfor{3}

\begin{enumerate}
\def\labelenumi{\arabic{enumi}.}
\item
  设 \(G\) 是一个李群，\(G_1\) 和 \(G_2\) 是 \(G\) 的李子群。假设 \(G_2\) 在 \(G\) 中具有有限余维，且 \(K\) 的特征为 0。则 \(G_1 \cap G_2\) 是 \(G\) 的李子群，其李代数为 \(L(G_1) \cap L(G_2)\)（仿照命题 29 及其推论 2 论证）。
\item
  假设 \(\mathbf{K}\) 是局部紧的。令 \(\mathbf{G}\) 是维数为 \(n\) 的李群。
\end{enumerate}

\begin{enumerate}
\def\labelenumi{(\alph{enumi})}
\tightlist
\item
  设 \(\phi\) 是 G 的具有有限核的 étale 端同态。令 \(\mu\) 是 G 的左 Haar 测度。则 \(\phi\) 是适当的，且
\end{enumerate}

\[
\phi (\mu) = \alpha \cdot \mu | \phi (G) \quad \text { where } \quad \alpha = \frac {\operatorname{Card} (\operatorname{Ker} \phi)}{\operatorname{mod} \det L (\phi)}.
\]

(使用命题 55。)
% semantic-fragments: legacy-input-seam
\relax{}
\begin{enumerate}
\def\labelenumi{(\alph{enumi})}
\setcounter{enumi}{1}
\tightlist
\item
  假设 G 是紧的. 设 \(\phi\) 是 G 的 étale 自同态. 则 \(\operatorname{Ker} \phi\) 是有限的, \(\phi(G)\) 在 G 中的指数有限, 并且
\end{enumerate}

\[
\frac {\operatorname{Card} (\operatorname{Ker} \phi)}{\operatorname{Card} (\mathrm{G} / \phi (\mathrm{G}))} = \operatorname{mod} \det \mathrm{L} (\phi).
\]

(利用 (a) 变换 \(\mu (\mathbf{G})\) .)

\begin{enumerate}
\def\labelenumi{(\alph{enumi})}
\setcounter{enumi}{2}
\tightlist
\item
  假设 G 是紧的且交换的. 设 \(r \in \mathbf{Z}\), 使得 r, 1 在 K 中均不为 0, 并设 \(r, 1 \neq 0\)\(\phi\) 是 G 上的自同态 \(x \mapsto x^r\). 则
\end{enumerate}

\[
\frac {\operatorname{Card} (\operatorname{Ker} \phi)}{\operatorname{Card} (\mathrm{G} / \phi (\mathrm{G}))} = (\mathrm{mod} r. 1) ^ {n}.
\]

(注意由 §2 命题 7 的推论, \(\mathrm{L}(\phi)\) 是比为 r 的伸缩变换.)\(r\)

\begin{enumerate}
\def\labelenumi{\arabic{enumi}.}
\setcounter{enumi}{2}
\tightlist
\item
  设 E 是由 2 行 2 列复对称矩阵组成的复向量空间. 对所有 \(s \in \mathbf{SL}(2, \mathbf{C})\), 定义 E 的自同构 \(\varphi(s)\), 使 \(\varphi(s).m = s.m.^t s\) .
\end{enumerate}

\begin{enumerate}
\def\labelenumi{(\alph{enumi})}
\tightlist
\item
  证明 \(\rho\) 是 \(\mathbf{SL}(2, \mathbf{C})\) 的解析线性表示. 将 \(\begin{pmatrix} a & b \\ b & c \end{pmatrix} \in \mathrm{E}\) 识别为 \((a, b, c) \in \mathbf{C}^3\) 后, \(\rho\begin{pmatrix} \alpha & \beta \\ \gamma & \delta \end{pmatrix}\) 的矩阵为
\end{enumerate}

\[
\left( \begin{array}{c c c} \alpha^ {2} & 2 \alpha \beta & \beta^ {2} \\ \alpha \gamma & \alpha \delta + \beta \gamma & \beta \delta \\ \gamma^ {2} & 2 \gamma \delta & \delta^ {2} \end{array} \right).
\]

\begin{enumerate}
\def\labelenumi{(\alph{enumi})}
\setcounter{enumi}{1}
\tightlist
\item
  证明 \(m \mapsto \det(m)\) 是非退化二次型
\end{enumerate}

\[
q (a, b, c) = a c - b ^ {2}
\]

在 E 上, 并且 \(\rho\) 是从 \(\mathbf{SL}(2,\mathbf{C})\) 到 \(\mathbf{SO}(q)\) 的满射态射, 核为 \(\{I, - I\}\) . (关于满射性, 注意 \(\mathbf{SO}(q)\) 中一个元素的矩阵的第 1 列可表示为 \((\alpha^2,\alpha \gamma ,\gamma^2)\) 的形式, 第 3 列可表示为 \((\beta^2,\beta \delta ,\delta^2)\) 的形式, 其中 \(\alpha^2\delta^2 +\beta^2\gamma^2 -2\alpha \gamma \beta \delta = 1\) . 然后注意, 在一个非各向同性平面上相同的 \(\mathbf{SO}(q)\) 中两个元素彼此相等.)

\begin{enumerate}
\def\labelenumi{(\alph{enumi})}
\setcounter{enumi}{2}
\tightlist
\item
  推出复李群 \(\mathbf{SO}(3, \mathbf{C})\) 同构于复李群 \(\mathbf{SL}(2, \mathbf{C}) / \{I, -I\}\) .
\end{enumerate}

\begin{enumerate}
\def\labelenumi{\arabic{enumi}.}
\setcounter{enumi}{3}
\tightlist
\item
  \begin{enumerate}
  \def\labelenumii{(\alph{enumii})}
  \tightlist
  \item
    设 \(q\) 是 \(\mathbf{R}^3\) 上符号为 (1, 2) 的二次型. 令 \(\mathbf{P}\) 为满足  的 \(m \in \mathbf{R}^3\) 的集合. 则 \(q(m) > 0\)\(\mathbf{P}\) 是两个不交凸锥 \(\mathbf{C}\) 和 \(-\mathbf{C}\) 的并. 若 \(s \in \mathbf{SO}(q)\), 则或者 \(s(\mathbf{C}) = \mathbf{C}\) 且 \(s(-\mathbf{C}) = -\mathbf{C}\), 或者 \(s(\mathbf{C}) = -\mathbf{C}\) 且 \(s(-\mathbf{C}) = \mathbf{C}\). 令 \(\mathbf{SO}^{+}(q)\) 为满足  的 \(s \in \mathbf{SO}(q)\) 的集合. 则 \(s(\mathbf{C}) = \mathbf{C}\)\(\mathbf{SO}^{+}(q)\) 是 \(\mathbf{SO}(q)\) 的开正规子群, 且在 \(\mathbf{SO}(q)\) 中的指数为 2 .
  \end{enumerate}
\end{enumerate}

\begin{enumerate}
\def\labelenumi{(\alph{enumi})}
\setcounter{enumi}{1}
\item
  仿照练习 3 的方法, 定义一个从 \(\mathbf{SL}(2,\mathbf{R})\) 到 \(\mathbf{SO}^{+}(q)\) 的满射态射, 其核为 \(\{I, - I\}\) . 推出实李群 \(\mathbf{SO}^{+}(q)\) 同构于实李群 \(\mathbf{SL}(2,\mathbf{R}) / \{I, - I\}\) .
\item
  对于 \((\xi, \eta) \in \mathbf{C}^2\) 和 \((\xi', \eta') \in \mathbf{C}^2\), 令 \(f((\xi, \eta), (\xi', \eta')) = \xi \bar{\xi}' - \eta \bar{\eta}'\) . 证明由  定义的 \(\mathbf{SL}(2, \mathbf{C})\) 的内自同构将 \(\frac{1}{\sqrt{2}} \begin{pmatrix} 1 & -i \\ -i & 1 \end{pmatrix}\)\(\mathbf{SU}(f)\) 映到 \(\mathbf{SL}(2, \mathbf{R})\) .
\end{enumerate}

\begin{enumerate}
\def\labelenumi{\arabic{enumi}.}
\setcounter{enumi}{4}
\tightlist
\item
  设 E 是由迹为零的 2 行 2 列厄米复矩阵组成的实向量空间. 对所有 \(s \in \mathbf{SU}(2, \mathbf{C})\), 定义 E 的自同构 \(\rho(s)\), 使 \(\rho(s)(m) = sms^*\) . 将元素
\end{enumerate}

\[
\left( \begin{array}{c c} x & y + i z \\ y - i z & - x \end{array} \right)
\]

与  中的元素 \((x,y,z)\) 对应. 仿照练习 3 的方法, 证明 \(R^{3}\)\(\rho\) 是从实李群 \(\mathbf{SU}(2,\mathbf{C})\) 到实李群 \(\mathbf{SO}(3,\mathbf{R})\) 的满射态射, 核为 \(\{I,-I\}\), 并且 \(\mathbf{SO}(3,\mathbf{R})\) 同构于 \(\mathbf{SU}(2,\mathbf{C})/\{I,-I\}\)

¶ 6. 设 F 是由 2 行 2 列厄米复矩阵组成的向量空间. 对所有 \(s \in \mathbf{SL}(2, \mathbf{C})\), 定义 F 的自同构 \(\sigma(s)\), 使 \(\sigma(s)(m) = sms^*\) . 将 F 中的元素 \(\begin{pmatrix} t + x & y + iz \\ y - iz & t - x \end{pmatrix}\) 与  中的元素 \((t, x, y, z)\) 对应, 并记\(\mathbf{R}^4\)

\[
q (t, x, y, z) = t ^ {2} - x ^ {2} - y ^ {2} - z ^ {2}.
\]

令 C 为满足 \((t, x, y, z) \in \mathbf{R}^+\)\(t^2 - x^2 - y^2 - z^2 > 0\), \(t > 0\) 的  的集合. 令 \(\mathbf{SO}^+(q)\) 为满足  的 \(g \in \mathbf{SO}(q)\) 的集合; 它是 \(g(\mathbf{C}) = \mathbf{C}\)\(\mathbf{SO}(q)\) 的开正规子群, 在 \(\mathbf{SO}(q)\) 中的指数为 2 (参见练习 4). 证明 \(\sigma\) 是从 \(\mathbf{SL}(2, \mathbf{C})\) 的底层实李群到实李群 \(\mathbf{SO}^+(q)\) 的满射态射, 核为 \(\{I, -I\}\) . (关于满射性, 使用练习 5.) 推出实李群 \(\mathbf{SO}^+(q)\) 同构于 \(\mathbf{SL}(2, \mathbf{C}) / \{I, -I\}\) 的底层实李群, 即 (练习 3) \(\mathbf{SO}(3, \mathbf{C})\) 的底层实李群 .

\begin{enumerate}
\def\labelenumi{\arabic{enumi}.}
\setcounter{enumi}{6}
\tightlist
\item
  \begin{enumerate}
  \def\labelenumii{(\alph{enumii})}
  \tightlist
  \item
    令 G 为范数为 1 的四元数集合. 这是实李群 \(\mathbf{H}^*\) 的李子群, 同胚于 \(\mathbf{S}_3\), 因而单连通, 参见《一般拓扑》第 XI 章.
  \end{enumerate}
\end{enumerate}

\begin{enumerate}
\def\labelenumi{(\alph{enumi})}
\setcounter{enumi}{1}
\tightlist
\item
  设 E 是纯四元数组成的实向量空间, 通过下式与 \(\mathbf{R}^3\) 对应
\end{enumerate}

\[
(x, y, z) \mapsto x i + y j + z k.
\]

对所有 \(g \in \mathbf{G}\), 定义  的自同构 \(\rho(g)\), 使 \(\mathbf{E}\)\(\rho(g)q = ggg^{-1}\) . 证明 \(\rho\) 是从李群 \(\mathbf{G}\) 到实李群 \(\mathbf{SO}(3, \mathbf{R})\) 的满射态射, 核为 \(\{1, -1\}\) . 推出李群 \(\mathbf{SO}(3, \mathbf{R})\) 同构于李群 \(G / \{1, -1\}\) .

\begin{enumerate}
\def\labelenumi{(\alph{enumi})}
\setcounter{enumi}{2}
\item
  将 \(\mathbf{C}\) 与  的子域 \(\mathbf{R} + \mathbf{R}i\) 对应. 映射 \(\mathbf{H}\)\((\lambda, q) \mapsto q\lambda\) 将 \(\mathbf{C} \times \mathbf{H}\) 映到 \(\mathbf{H}\), 使 \(\mathbf{H}\) 成为复向量空间; 通过选取基 \(\mathbf{C}\), 它与 \(\mathbf{C}^2\) 对应. 对所有 \((1, j)\)\(g \in \mathbf{G}\), 令 \(\sigma(g)\) 为由 \(\sigma(g)q = gq\) 定义的该向量空间的自同构. 证明 \(\sigma\) 是从实李群 G 到实李群 \(\mathbf{SU}(2, \mathbf{C})\) 的同构. 因而由 \((a)\), 后者是单连通的.
\item
  对于 \((q_{1}, q_{2}) \in \mathbf{G} \times \mathbf{G}\), 令 \(\tau(q_{1}, q_{2})\) 为从  到自身的映射 \(q \mapsto q_{1}q\bar{q}_{2}\). 证明 \(\mathbf{H} = \mathbf{R}^{4}\)\(\tau\) 是从实李群 \(\mathrm{G} \times \mathrm{G}\) 到实李群 \(\mathbf{SO}(4, \mathbf{R})\) 的满射态射, 核为 \(\mathrm{N} = \{(1, 1), (-1, -1)\}\), 因而 \(\mathbf{SO}(4, \mathbf{R})\) 同构于 \((\mathrm{G} \times \mathrm{G}) / \mathrm{N}\) .
\end{enumerate}

\begin{enumerate}
\def\labelenumi{\arabic{enumi}.}
\setcounter{enumi}{7}
\tightlist
\item
  设 \(\mathrm{G}\) 是有限维连通实李群, \(\mathbf{M}\) 是有限维连通的 \(\mathrm{C}^{\infty}\) 类流形; 给定 \(\mathrm{C}^{\infty}\)\(\mathrm{G}\) 在 \(\mathbf{M}\) 上的  类右作用法则. 对所有 \(x \in \mathbf{M}\), 令 \(\varphi(x)\) 为相应的轨道映射.
\end{enumerate}

\begin{enumerate}
\def\labelenumi{(\alph{enumi})}
\item
  G 在 M 上传递作用, 当且仅当对所有 \(x \in M\), 从 L(G) 到  的映射 \(\mathrm{T}_{e}(\varphi(x))\) 是满射. (要看条件的充分性, 注意此时 M 中的每个 G-轨道都是开集.)\(\mathrm{T}_{x}(\mathrm{M})\)
\item
  假设 G 在 M 上传递作用. 我们把 M 识别为齐性空间 H\textbar G, 其中 H 是 M 中一点 \(x_{0}\) 的稳定子. 若存在 M 的一个 \(C^{\omega}\) 类闭子流形 V, 满足 \(0 < dim V < dim M\), 且对每个变换 V.s (其中 \(s \in G\)) 而言, V.s 要么等于 V, 要么与 V 不交, 则称 G 在 M 上的作用是非本原的. 这当且仅当存在 G 的李子群 L, 使得 \(H \subset L \subset G\) 且 \(\dim(H) < \dim(L) < \dim(G)\). 若不存在这样的 L, 则称 G 在 M 上的作用是本原的; 只要 L(G) 与 L(H) 之间不存在不同于 L(G) 和 L(H) 的子代数, 这就足够了.
\item
  在 (b) 的假设下, 令 \(\mathcal{L}\) 为 L(G) 在与 G 的作用相联系的无穷小作用法则下的像. 令 \(\mathcal{E}\) 为 M 上 \(C^\infty\) 类函数的代数, m 为 \(\mathcal{E}\) 的极大理想, 由 \(\mathcal{E}\) 中在 \(x_0\) 处为零的函数组成; 对于 , 记 \(\mathcal{L}_p\) 为满足对所有 \(p = -1, 0, 1, 2, \ldots\) 都有  的 \(x \in \mathcal{L}\) 的集合; 于是 \(\theta(\mathbf{X}) f \in m^{p+1}\)\(f \in \mathcal{E}\)\(\mathcal{L}_{-1} = \mathcal{L}\), 而 \(\mathcal{L}_0\) 是 L(H) 在无穷小作用法则下的像. 若 (U, \(\phi, \mathbf{R}^n\) ) 是以 \(x_0\) 为中心的 M 的图, 且 \(\mathbf{X}_1, \ldots, \mathbf{X}_n\) 是由该图定义的 U 上向量场, 则 \(\mathcal{L}_p\) 的元素是形如 \(a_1\mathbf{X}_1 + \cdots + a_n\mathbf{X}_n\) 的元素, 其中 \(a_1, \ldots, a_n\) 属于 \(m^p|\text{U}\). 证明 \([\mathcal{L}_p, \mathcal{L}_q] \subset \mathcal{L}_{p,q}\) (约定 \(\mathcal{L}_{-2} = \mathcal{L}\)). 若存在 Y \(\in \mathcal{L}_p\) (其中 \(p \geqslant 0\)) 使得 Y \(\notin \mathcal{L}_{p+1}\), 则存在 X \(\in\) L 使得 {[}Y, X{]} \(\notin \mathcal{L}_p\). \(\mathcal{L}_p\) 是 \(\mathcal{L}\) 的李子代数. 对于 \(p \geqslant 0\), \(\mathcal{L}_p\) 是 \(\mathcal{L}_0\) 的理想. 令 \(\rho\) 为 H 在 T \(_{x_0}\) (M) 上的典范线性表示. H/(Ker \(\rho\) ) 的李代数同构于 \(\mathcal{L}_0/\mathcal{L}_1\).
\item
  进一步假设各 \(\mathcal{L}_p\) 的交为 \(\{0\}\). 则存在最大的指标 r, 使 \(r\)\(\mathcal{L}_r \neq \{0\}\), 且所有指标  的 \(\mathcal{L}_f\) 两两不同. 对于 \(\leqslant r\)\(p \geqslant 0\),
\end{enumerate}

\[
\dim (\mathcal {L} _ {p} / \mathcal {L} _ {p + 1}) \leqslant \binom {n + p} {p + 1}.
\]

\begin{enumerate}
\def\labelenumi{(\alph{enumi})}
\setcounter{enumi}{4}
\tightlist
\item
  若 M 是解析流形且 G 在 M 上解析地作用, 则满足 (d) 的假设.
\end{enumerate}

\begin{enumerate}
\def\labelenumi{\arabic{enumi}.}
\setcounter{enumi}{8}
\tightlist
\item
  在命题 30 的记号下, HH' 可以是 G 中稠密而不同于 G 的开子集 (取 G = GL(2, R), H 为上三角子群, H' 为严格下三角子群).
\end{enumerate}

